% Part: methods
% Chapter: proofs
% Section: cant-do-it

\documentclass[../../../include/open-logic-section]{subfiles}

\begin{document}

\olfileid{mth}{prf}{cnt}

\olsection{زه ئې نۀ شم کولے!{}}

موږ ټول داسې پړاو ته رسېږو چې د پرېښودلو زړه مو کېږي۔
خو تاسو دا کار \emph{کولے شئ}۔ استاد، د تدريس مرستيال
او ستاسو ملګري زده کوونکي مرسته کولے شي۔ مرسته ترې
وغواړئ!{} دلته څو مشورې دي چې له سخت حالت څخه مو وساتي،
او دا هم ښيي چې د پرېښودلو د احساس پر وخت څه وکړئ۔

د دې دپاره چې مسئلې په برياليتوب حل کړے شئ، دا کارونه وکړئ:
\begin{enumerate}
\item \emph{څومره چې کېدلے شي، هومره وختي پيل وکړئ۔}
  په ټول سمسټر کښې بوختېږو او زموږ ډېرو ته د کار ځنډول
  ستونزه وي؛ يو له غوره کارونو دا دے چې کورنۍ دندې
  وختي پيل کړئ۔ په دې توګه، که ايسار شئ، د حل د لټولو
  وخت به لرئ (چې ژړا نۀ وي)۔
\item \emph{له خپلو ټولګيوالو سره خبرې وکړئ}۔ تاسو يوازې
  نۀ يئ۔ ښايي په ټولګي کښې نورو ته هم ستونزې وي، خو
  کېدے شي د هغوى ستونزې په نورو څيزونو کښې وي۔ له
  ملګرو سره پرې خبرې کول د مسئلې بله زاويه درښودلے شي
  چې د حل لاره پرانيزي۔ البته، يوازې د هغوى حل مه نقلوئ:
  اشاره ترې وغواړئ، يا ووايئ چې چېرته ايسار يئ او د بل
  ګام پوښتنه وکړئ۔ او کله چې تاسو پرې پوه شئ، بيا د
  هغوى مرسته وکړئ۔ د بل چا مرسته او د مطلب تشريح کول
  ستاسو خپله پوهه هم ښه کوي۔
\item \emph{مرسته وغواړئ۔} ډېرې وسيلې درته شته؛ استاد
  او د تدريس مرستيال ستاسو دپاره حاضر دي او ستاسو
  برياليتوب \emph{غواړي}۔ هغوى بايد د مسئلې په حلولو او
  د بهير د هغه ځاے په پېژندلو کښې مرسته وکړے شي چې
  تاسو پکښې ستونزه لرئ۔
\item \emph{دمه وکړئ۔} که ايسار يئ، \emph{کېدے شي}
  وجه ئې دا وي چې له ډېرې مودې راهيسې هماغې مسئلې ته
  ګورئ۔ لنډه دمه وکړئ، يوه پياله چاے وڅښئ، يا يو څه
  وخت په بله مسئله کار وکړئ؛ بيا په تازه ذهن هغې ته
  راوګرځئ۔ يوه شپه ئې پرېږدئ او تر خوب وروسته بيا پرې
  فکر وکړئ۔
\end{enumerate}

ګورئ چې دا لارې له تاسو غواړي ثبوت له مخکې نه وختي
پيل کړے وي؟ که د سپارلو له نېټې يوه ورځ وړاندې د شپې
په دوو بجو مو ثبوت پيل کړے وي، ښايي دا لارې دومره
مرستندويې نۀ وي۔

ښايي دا خبرې ډېرې ناامېدوونکې ښکاره شي، خو د ثبوت
حلول داسې ننګونه ده چې په پاې کښې ګټه لري۔ ځينې
خلک همدا کار د مسلک په توګه کوي؛ نو ضرور به پکښې
د خوند يو څه وي۔ د نږدې هر کار په شان د مسئلو
حلول او د ثبوتونو کول تمرين غواړي۔ ښايي داسې
ټولګيوال ووينئ چې دا کار ورته اسان وي؛ غالباً
هغوى له مخکې ډېر تمرين کړے دے۔ هڅه وکړئ چې ډېر
په آسانۍ لاس ترې وانۀ خلئ۔

که په کومه ځانګړې مسئله وخت يا زغم درنه ختم شي،
پروا نۀ کوي۔ دا معنا نۀ لري چې تاسو ناپوه يئ يا
هېڅکله به پرې پوه نۀ شئ۔ له استاد يا بل زده
کوونکي معلومه کړئ چې څنګه حلېږي، او دا وپېژنئ
چې چېرته مو تېروتنه کړې يا ايسار شوي يئ، څو
راتلونکے ځل له ورته ستونزې سره د مخ کېدو پر
وخت ترې ځان وساتئ۔ بيا هڅه وکړئ چې حل ته له
کتلو پرته ئې پخپله وکړئ۔ او راتلونکے ځل وختي
پيل وکړئ او وختي مرسته وغواړئ۔

\end{document}
